% Lösungen Einheit 1 von 4 – Zuordnungen darstellen (zusammenbau v0.1)
\einheitenkopf{Einheit 1 von 4 · Zuordnungen darstellen}

\erg{8}{a) (1|15), (2|30), (3|45) \quad b) (1|2), (2|4), (3|6), (4|8) \quad c) 6 km; 4 h \quad d) (1|20), (2|40), (3|60), (4|80) \quad e) Nein – halbe Kinokarten gibt es nicht; nur die Punkte selbst haben einen Sinn. \quad f) steigende Gerade durch den Ursprung \quad g) Mia hat die Achsen vertauscht: richtig sind (1|3) und (2|6), nicht (3|1) und (6|2). \quad h) (0|70), (1|95), (2|135), (3|185), (4|260)}
\erg{9}{a) 5 (20 : 4 = 5) \quad b) 1 Kästchen = 5 km, denn 60 km sind dann 12 Kästchen. \quad c) zum Beispiel 1 Kästchen = 1 Heft, Achse von 0 bis 6 beschriftet \quad d) 18 Kästchen; nein, die Achse hat nur 16. \quad e) 40 m (800 : 20 = 40) \quad f) waagerecht: Zeit in min; senkrecht: Strecke in m \quad g) fehlende Werte 90 l und 30 l; zum Beispiel 1 Kästchen = 2 min (30 min = 15 Kästchen) und 1 Kästchen = 5 l (90 l = 18 Kästchen); Punkte (0|90), (5|80), (10|70), (15|60), (20|50), (25|40), (30|30)}
\erg{10}{a) y = 3 · x}
\erg{11}{a) Lea hat Kästchen gezählt statt Werte abgelesen: 5 Kästchen sind 5 · 2 km = 10 km. \quad b) Nein – zu jeder Uhrzeit gehört genau eine Temperatur; sonst wäre es keine Zuordnung. \quad c) (1|1,5), (2|3), (3|4,5) \quad d) um 14 Uhr; 11 °C wärmer}
